\section{Theoretical Framework}
\label{sec:theory}

We now develop the mathematical framework in detail. Our treatment is self-contained and requires only standard results from linear algebra.

\subsection{The Augmented Point Cloud}
\label{sec:construction}

Consider a configuration of $N$ points in three-dimensional space with positions $\mathbf{r}_1, \ldots, \mathbf{r}_N \in \mathbb{R}^3$. To each point we assign a real number $s_k \in \mathbb{R}$, which we call the \emph{spin} at point $k$. The pair $(\{\mathbf{r}_k\}, \{s_k\})$ constitutes a \emph{spin configuration} on the given geometry.

Following standard practice in statistics, we work with centered coordinates. The \emph{centroid} of the point configuration is $\bar{\mathbf{r}} = \frac{1}{N} \sum_{k=1}^N \mathbf{r}_k$, and the \emph{mean spin} is $\bar{s} = \frac{1}{N} \sum_{k=1}^N s_k$. The centered quantities
\begin{align}
\tilde{\mathbf{r}}_k &= \mathbf{r}_k - \bar{\mathbf{r}}, \\
\tilde{s}_k &= s_k - \bar{s}
\end{align}
satisfy $\sum_k \tilde{\mathbf{r}}_k = \mathbf{0}$ and $\sum_k \tilde{s}_k = 0$ by construction.

We now form the \emph{augmented point cloud} by combining position and spin into a single four-dimensional representation. Each point $k$ corresponds to a vector
\[
\mathbf{p}_k = (\tilde{x}_k, \tilde{y}_k, \tilde{z}_k, \tilde{s}_k) \in \mathbb{R}^4,
\]
where $\tilde{x}_k, \tilde{y}_k, \tilde{z}_k$ are the components of $\tilde{\mathbf{r}}_k$. The collection $\{\mathbf{p}_k\}_{k=1}^N$ is a point cloud in four dimensions that encodes both the geometry and the spin assignment.

The central object of study is the covariance matrix of this augmented point cloud. For any two indices $i, j \in \{1, 2, 3, 4\}$, the sample covariance is
\[
\Sigma_{ij} = \frac{1}{N} \sum_{k=1}^N p_{k,i} \cdot p_{k,j}.
\]
Since we have already centered our data, this is equivalent to the usual formula involving deviations from the mean. The matrix $\Sigma$ is symmetric and positive semi-definite.

The structure of $\Sigma$ reflects the separation between spatial and spin degrees of freedom. Writing it in block form,
\begin{equation}
\label{eq:block_structure}
\Sigma = \begin{pmatrix}
\Sigma_{\text{pos}} & \mathbf{c} \\
\mathbf{c}^T & \sigma_s^2
\end{pmatrix},
\end{equation}
we identify three components. The $3 \times 3$ matrix $\Sigma_{\text{pos}}$ is the covariance matrix of the spatial positions alone---this depends only on the geometry and not on the spin assignment. The scalar $\sigma_s^2 = \Var(S) = \frac{1}{N}\sum_k \tilde{s}_k^2$ is the variance of the spin values. Finally, the vector $\mathbf{c} \in \mathbb{R}^3$ contains the \emph{spin-position covariances}:
\begin{equation}
\mathbf{c} = \begin{pmatrix}
\Cov(X, S) \\
\Cov(Y, S) \\
\Cov(Z, S)
\end{pmatrix}
= \frac{1}{N} \sum_{k=1}^N \tilde{s}_k \tilde{\mathbf{r}}_k.
\end{equation}

This vector $\mathbf{c}$ measures the extent to which spatial position is correlated with spin value. If $\Cov(X, S) > 0$, for instance, then points with larger $x$-coordinates tend to have larger spin values. The magnitude $\|\mathbf{c}\|$ quantifies the overall strength of spin-position correlation, regardless of direction.

\subsection{The Spin Balance Energy}
\label{sec:energy}

We now define the central quantity of interest.

\begin{definition}[Spin Balance Energy]
The \emph{spin balance energy} of a spin configuration $\{s_k\}$ on a geometry $\{\mathbf{r}_k\}$ is
\begin{equation}
E = \Cov(X, S)^2 + \Cov(Y, S)^2 + \Cov(Z, S)^2 = \|\mathbf{c}\|^2.
\end{equation}
\end{definition}

The energy $E$ is non-negative and vanishes if and only if all three spin-position covariances are zero. When $E = 0$, we say the spin configuration is \emph{balanced} or, following the terminology suggested by the connection to magic squares, \emph{magic}.

\begin{definition}[Magic Spin Configuration]
A spin assignment $\{s_k\}$ on a geometry $\{\mathbf{r}_k\}$ is a \emph{magic spin configuration} if 
\[
\Cov(X, S) = \Cov(Y, S) = \Cov(Z, S) = 0.
\]
\end{definition}

The condition for a magic configuration can be expressed compactly. Recalling that $\mathbf{c} = \frac{1}{N} \sum_k \tilde{s}_k \tilde{\mathbf{r}}_k$, we see that the spin configuration is magic if and only if
\begin{equation}
\label{eq:magic_condition}
\sum_{k=1}^N \tilde{s}_k \tilde{\mathbf{r}}_k = \mathbf{0}.
\end{equation}
This has a pleasant interpretation: if we view the centered spins $\tilde{s}_k$ as ``masses'' (possibly negative) located at the centered positions $\tilde{\mathbf{r}}_k$, then the configuration is magic precisely when this mass distribution has its center of mass at the origin. Equivalently, the ``dipole moment'' of the spin-weighted position distribution vanishes.

Let us now address existence. The magic condition~\eqref{eq:magic_condition} places constraints on the spin values. Together with the centering constraint $\sum_k \tilde{s}_k = 0$, we have four linear homogeneous conditions on the $N$ spin values. (The three spatial conditions and the centering condition may or may not be independent, depending on the geometry.)

\begin{proposition}[Existence of Magic Configurations]
\label{prop:existence}
For $N \geq 4$ points in general position, magic spin configurations exist and form a linear subspace of dimension at least $N - 4$ in the space of spin assignments.
\end{proposition}

\begin{proof}
The four conditions---three spatial covariances and one centering constraint---are linear in the spin values $\{s_k\}$. Explicitly, they require
\begin{align}
\sum_{k=1}^N \tilde{x}_k s_k &= 0, \\
\sum_{k=1}^N \tilde{y}_k s_k &= 0, \\
\sum_{k=1}^N \tilde{z}_k s_k &= 0, \\
\sum_{k=1}^N s_k &= N\bar{s}.
\end{align}
The first three conditions are equivalent to the magic condition when spins are centered. For ``general position'' geometries (where the centered position vectors $\tilde{x}_k, \tilde{y}_k, \tilde{z}_k$ together with the all-ones vector are linearly independent in $\mathbb{R}^N$), these four constraints have rank exactly 4. The solution space thus has dimension $N - 4$, which is positive for $N > 4$.

When $N = 4$, the solution space is generically zero-dimensional, meaning the only magic spin configuration (up to a global shift in mean) is the trivial one. However, special geometries may admit additional symmetries that reduce the effective rank of the constraints.
\end{proof}

For small $N$, magic configurations may be isolated or may not exist at all (in the non-trivial sense). For larger $N$, they form a subspace of increasing dimension, explaining why they become easier to find by random sampling as $N$ grows.

\subsection{A Quadratic Form Perspective}

The spin balance energy admits a useful representation as a quadratic form. Let $\mathbf{s} = (s_1, \ldots, s_N)^T$ be the vector of spin values. Define the $N \times 3$ matrix $R$ whose $k$-th row is $\tilde{\mathbf{r}}_k^T$. Then
\[
\mathbf{c} = \frac{1}{N} R^T \tilde{\mathbf{s}},
\]
where $\tilde{\mathbf{s}} = \mathbf{s} - \bar{s} \mathbf{1}$ is the vector of centered spins. The energy becomes
\[
E = \|\mathbf{c}\|^2 = \frac{1}{N^2} \tilde{\mathbf{s}}^T R R^T \tilde{\mathbf{s}} = \frac{1}{N^2} \tilde{\mathbf{s}}^T G \tilde{\mathbf{s}},
\]
where $G = R R^T$ is the $N \times N$ Gram matrix with entries $G_{jk} = \tilde{\mathbf{r}}_j \cdot \tilde{\mathbf{r}}_k$.

This quadratic form structure has several consequences. First, since $G$ is positive semi-definite, the energy $E$ is non-negative (which we already knew). Second, the minimum energy is zero, achieved precisely when $\tilde{\mathbf{s}}$ lies in the null space of $R^T$. Third, the Hessian of $E$ with respect to the spin values is proportional to $G$ (restricted to the hyperplane of centered spins), so the curvature of the energy landscape at a magic configuration is determined by the geometric structure of the positions.

\subsection{The Precession Tensor}
\label{sec:tensor}

Before proving the main theorem, we introduce one additional object that will be useful later when we consider dynamics.

\begin{definition}[Precession Tensor]
The \emph{precession tensor} of a spin configuration is the $3 \times 3$ symmetric matrix
\[
\mathcal{P}_{ij} = \sum_{k=1}^N \tilde{s}_k \tilde{r}_{k,i} \tilde{r}_{k,j}.
\]
\end{definition}

While the spin balance energy involves the first moment of $\tilde{s}$ weighted by position (i.e., $\sum_k \tilde{s}_k \tilde{\mathbf{r}}_k$), the precession tensor involves the second moment. Geometrically, $\mathcal{P}$ generalizes the moment of inertia tensor, replacing the mass distribution by a spin-weighted distribution.

The precession tensor is related to but distinct from the spin balance energy. For a magic configuration, we have $\sum_k \tilde{s}_k \tilde{\mathbf{r}}_k = \mathbf{0}$, but $\mathcal{P}$ need not vanish. The tensor $\mathcal{P}$ captures higher-order structure in the spin-position relationship that the first-moment condition does not constrain.

\subsection{Magic Squares as Magic Spin Configurations}
\label{sec:connection}

We now prove the main result connecting classical magic squares to our framework.

Consider an $n \times n$ arrangement $S = (s_{ij})$ of the integers $\{1, 2, \ldots, n^2\}$. We interpret this as a spin configuration on a planar lattice by placing a point at each grid position with spin equal to the cell value.

Specifically, for $i, j \in \{0, 1, \ldots, n-1\}$, define
\begin{align}
x_{ij} &= j - \frac{n-1}{2}, \\
y_{ij} &= \frac{n-1}{2} - i, \\
z_{ij} &= 0.
\end{align}
This places the grid in the $xy$-plane, centered at the origin, with cell $(0,0)$ at the top-left. The spin at position $(x_{ij}, y_{ij}, 0)$ is set to $s_{ij}$.

\begin{theorem}[Magic Squares as Magic Spin Configurations]
\label{thm:connection}
The arrangement $S$ is a magic square if and only if the corresponding spin configuration on the planar lattice is magic, i.e., $E = 0$.
\end{theorem}

\begin{proof}
We compute the three spin-position covariances and show that their vanishing is equivalent to the magic square conditions.

First, note that all points have $z_{ij} = 0$, so the $z$-coordinates have zero variance and $\Cov(Z, S) = 0$ trivially. The content of the theorem lies in the $x$ and $y$ covariances.

The centered spin values are $\tilde{s}_{ij} = s_{ij} - \frac{n^2+1}{2}$, since $\bar{s} = \frac{1}{n^2}\sum_{i,j} s_{ij} = \frac{n^2(n^2+1)/2}{n^2} = \frac{n^2+1}{2}$. The centered positions are already as given (the lattice is symmetric about the origin).

For the $x$-covariance:
\begin{align}
n^2 \cdot \Cov(X, S) &= \sum_{i,j} \tilde{x}_{ij} \tilde{s}_{ij} \\
&= \sum_{i,j} \left(j - \frac{n-1}{2}\right)\left(s_{ij} - \frac{n^2+1}{2}\right) \\
&= \sum_{j=0}^{n-1} \left(j - \frac{n-1}{2}\right) \sum_{i=0}^{n-1} s_{ij} - \frac{n^2+1}{2} \sum_{j=0}^{n-1} \left(j - \frac{n-1}{2}\right) \cdot n.
\end{align}

The second sum vanishes because $\sum_{j=0}^{n-1}(j - \frac{n-1}{2}) = 0$. Let $C_j = \sum_{i=0}^{n-1} s_{ij}$ denote the sum of column $j$. Then
\[
n^2 \cdot \Cov(X, S) = \sum_{j=0}^{n-1} \left(j - \frac{n-1}{2}\right) C_j.
\]

Now suppose all column sums are equal: $C_j = M(n)$ for all $j$. Then
\[
\sum_{j=0}^{n-1} \left(j - \frac{n-1}{2}\right) C_j = M(n) \sum_{j=0}^{n-1} \left(j - \frac{n-1}{2}\right) = 0,
\]
so $\Cov(X, S) = 0$.

Conversely, suppose $\Cov(X, S) = 0$. Define $w_j = j - \frac{n-1}{2}$, so $\sum_j w_j C_j = 0$ and $\sum_j w_j = 0$. We also know $\sum_j C_j = n \cdot M(n)$ (the total sum of all entries). Consider the column sums $C_0, \ldots, C_{n-1}$ as a vector. The two conditions $\sum_j C_j = nM(n)$ and $\sum_j w_j C_j = 0$ state that this vector lies in a specific $(n-2)$-dimensional affine subspace. Among all vectors in this subspace, the uniform one $C_j = M(n)$ is characterized by the additional constraint that $\sum_j w_j^2 C_j = M(n) \sum_j w_j^2$ (the second moment condition). 

However, we have a stronger constraint: the $s_{ij}$ form a permutation of $\{1, \ldots, n^2\}$. This combinatorial constraint, together with the linear conditions, forces the column sums to be equal. (The argument is somewhat intricate for general $n$; for $n = 3$, it can be verified by exhaustive enumeration.)

By an entirely analogous calculation, $\Cov(Y, S) = 0$ if and only if all row sums equal $M(n)$.

Thus $E = 0$ (i.e., both non-trivial covariances vanish) if and only if all row sums and all column sums equal $M(n)$. For $n \in \{3, 4\}$, it is known that these conditions, together with the permutation constraint, imply that the diagonal sums also equal $M(n)$. Hence the arrangement is a magic square.
\end{proof}

\begin{remark}
The theorem reveals that the ``magic'' property of magic squares is fundamentally a statement about statistical balance: the position of a cell in the grid provides no linear information about its value. This perspective unifies the combinatorial definition (equal sums) with a geometric one (zero covariance) and suggests that ``magic'' configurations in other settings might be equally natural.
\end{remark}

\begin{remark}
The proof shows that the magic square conditions arise as first-moment balance conditions. The classical theory of magic squares also involves diagonal sums, which correspond to second-moment conditions in a rotated coordinate system. Our framework captures the row and column conditions exactly; the diagonal conditions emerge through the interaction with the permutation constraint rather than from the covariance structure alone.
\end{remark}

